
\begingroup
\setlength{\textfloatsep}{0.6em}
\setlength{\floatsep}{0.6em}
\setlength{\intextsep}{0.6em}
\setlength{\abovecaptionskip}{0.35em}
\setlength{\belowcaptionskip}{0.2em}
\lstset{
  aboveskip=0.2em,
  belowskip=0.2em,
  frame=none,
  numbers=none,
  xleftmargin=0pt,
  framexleftmargin=0pt,
}

\chapter{API Reference and Sample Payloads}

This appendix summarizes the VietFlood API gateway endpoints and provides representative JSON responses. Token strings, user identifiers, and media URLs are masked with placeholders.

\section{Endpoint Inventory}

\begin{table}[H]
\centering
\caption{REST endpoints exposed by the VietFlood API gateway.}
\label{tab:appendix-api-endpoints}
\scriptsize
\setlength{\tabcolsep}{3pt}
\renewcommand{\arraystretch}{1.1}
\begin{tabularx}{\textwidth}{|l|p{4.6cm}|p{3.1cm}|X|}
\hline
\textbf{Method} & \textbf{Endpoint} & \textbf{Access} & \textbf{Description} \\
\hline
\multicolumn{4}{|l|}{\textbf{Authentication}} \\
\hline
POST & \url{/auth/register} & Public & Create a new user account. \\
\hline
POST & \url{/auth/sign_in} & Public & Authenticate and obtain access and refresh tokens. \\
\hline
POST & \url{/auth/refresh_token} & Refresh token & Issue a new access token and rotate the refresh token. \\
\hline
POST & \url{/auth/logout} & Public & Logout and revoke refresh token (if provided). \\
\hline
POST & \url{/auth/refresh} & Public & Refresh helper endpoint used by the gateway (payload passthrough). \\
\hline
GET & \url{/auth/profile} & JWT & Fetch the current user profile. \\
\hline
PUT & \url{/auth/update} & JWT & Update the current user's profile. \\
\hline
PUT & \url{/auth/update/user/:id} & JWT + Role (admin/relief) & Update a user profile by user id. \\
\hline
DELETE & \url{/auth/delete/:id} & JWT + Role (admin) & Delete a user by user id. \\
\hline
GET & \url{/auth/all} & JWT + Role (admin/relief) & List all users. \\
\hline
GET & \url{/auth/relief/users} & JWT + Role (admin/relief) & List users (alias of \url{/auth/all}). \\
\hline
GET & \url{/auth/relief/users/:id} & JWT + Role (admin/relief) & Fetch a user by user id. \\
\hline
\multicolumn{4}{|l|}{\textbf{Reports}} \\
\hline
POST & \url{/reports/create} & JWT & Create a new flood report (supports evidence file upload). \\
\hline
GET & \url{/reports} & JWT + Role (admin/relief) & Fetch all reports for admin/relief review. \\
\hline
GET & \url{/reports/user} & JWT + Role (citizen) & Fetch reports created by the current citizen user. \\
\hline
GET & \url{/reports/:id} & JWT + Role (admin/relief) & Fetch reports for a specific user id. \\
\hline
PUT & \url{/reports/update/:id} & JWT & Update a report owned by the current user (supports evidence upload). \\
\hline
PUT & \url{/reports/update/:id/admin/:userId} & JWT + Role (admin) & Update a report for a specific user id (admin workflow). \\
\hline
PUT & \url{/reports/relief/:id/user/:userId} & JWT + Role (admin/relief) & Update a report for a specific user id (relief workflow). \\
\hline
DELETE & \url{/reports/:id} & JWT & Delete a report owned by the current user. \\
\hline
DELETE & \url{/reports/update/:id/admin/:userId} & JWT + Role (admin) & Delete a report for a specific user id (admin workflow). \\
\hline
DELETE & \url{/reports/relief/:id/user/:userId} & JWT + Role (admin/relief) & Delete a report for a specific user id (relief workflow). \\
\hline
\end{tabularx}
\end{table}

\section{Authentication: Sign In}

\textbf{Endpoint:} \texttt{POST /auth/sign\_in}

\textbf{Response body:}
\begin{lstlisting}[language=JavaScript, caption=Sign In Response, label=lst:signin-response]
{
  "accessToken": "<JWT_ACCESS_TOKEN>",
  "refresh_token": "<JWT_REFRESH_TOKEN>"
}
\end{lstlisting}

\section{Authentication: Token Refresh}

\textbf{Endpoint:} \texttt{POST /auth/refresh\_token}

\textbf{Response body:}
\begin{lstlisting}[language=JavaScript, caption=Refresh Token Response, label=lst:refresh-token-response]
{
  "access_token": "<JWT_ACCESS_TOKEN>",
  "refresh_token": "<JWT_REFRESH_TOKEN>"
}
\end{lstlisting}

\section{User: Profile Payload}

\textbf{Endpoint:} \texttt{GET /auth/profile}

\textbf{Response body:}
\begin{lstlisting}[language=JavaScript, caption=User Profile Response, label=lst:profile-response]
{
  "id": 7,
  "email": "<user_email>",
  "username": "<username>",
  "phone": "<phone_number>",
  "role": "citizen",
  "first_name": "<first_name>",
  "middle_name": null,
  "last_name": "<last_name>",
  "date_of_birth": "2000-01-01",
  "address_line": "<address_line>",
  "ward": "<ward>",
  "province": "<province>",
  "created_at": "2026-05-17T00:00:00.000Z",
  "updated_at": "2026-05-17T00:00:00.000Z"
}
\end{lstlisting}

\section{Reports: Create Payload}

\textbf{Endpoint:} \texttt{POST /reports/create}

\textbf{Response body:}
\begin{lstlisting}[language=JavaScript, caption=Create Report Response, label=lst:create-report-response]
{
  "success": true,
  "message": "Create report successfully",
  "reportId": 41,
  "category": ["flood"],
  "evidences": [
    {
      "url": "<CLOUDINARY_URL>",
      "publicId": "<CLOUDINARY_PUBLIC_ID>",
      "resourceType": "image"
    }
  ],
  "images": [
    "<CLOUDINARY_URL>"
  ]
}
\end{lstlisting}

\chapter{User Interface Snapshots}

Key user interface screens are shown here. These images support the main report without repeating the implementation details.

\FloatBarrier
\section{Mobile and Dashboard Screens}

\begin{figure}[H]
    \centering
    \begin{minipage}[t]{0.30\linewidth}
        \centering
        \includegraphics[width=0.88\linewidth,keepaspectratio]{figure4-11-mobile-create-report.jpg}
        \caption{Mobile report creation screen (citizen workflow).}
        \label{fig:appendix-citizen-create-report}
    \end{minipage}
    \hfill
    \begin{minipage}[t]{0.30\linewidth}
        \centering
        \includegraphics[width=0.88\linewidth,keepaspectratio]{figure4-12-mobile-evidence-upload.jpg}
        \caption{Mobile evidence selection and upload screen.}
        \label{fig:appendix-mobile-evidence-upload}
    \end{minipage}
    \hfill
    \begin{minipage}[t]{0.30\linewidth}
        \centering
        \includegraphics[width=0.88\linewidth,keepaspectratio]{figure4-13-relief-report-list.jpg}
        \caption{Relief report list screen for browsing incoming records.}
        \label{fig:appendix-relief-report-list}
    \end{minipage}
    
    \vspace{0.55em}
    \begin{minipage}[t]{0.82\linewidth}
        \centering
        \includegraphics[width=\linewidth,keepaspectratio]{figure4-15-admin-dashboard-overview.jpg}
        \caption{Administrator dashboard overview with report list and filters.}
        \label{fig:appendix-admin-dashboard}
    \end{minipage}
\end{figure}
\FloatBarrier

\chapter{Prototype Validation Checklist}

This appendix lists system-level test cases used to validate the VietFlood prototype. The cases focus on end-to-end behavior through the API gateway, including authentication, role checks, report workflow, and evidence handling.

\section{Auth and User Scenarios}

\begin{table}[H]
\centering
\caption{Authentication and user management validation scenarios.}
\label{tab:appendix-auth-test-cases}
\renewcommand{\arraystretch}{1.15}
\scriptsize
\setlength{\tabcolsep}{3pt}
\begin{tabularx}{\textwidth}{|p{2.0cm}|p{3.8cm}|X|p{3.2cm}|}
\hline
\textbf{ID} & \textbf{Feature} & \textbf{Input / Steps (summary)} & \textbf{Expected Result} \\
\hline
AUTH-01 & Sign in success (\texttt{POST /auth/sign\_in}) & Provide valid \texttt{username} and \texttt{password}. & Response contains \texttt{accessToken} and \texttt{refresh\_token}. \\
\hline
AUTH-02 & Sign in failure & Provide valid username with incorrect password. & Request rejected with an authentication error (no token returned). \\
\hline
AUTH-03 & Access protected route without JWT & Call \texttt{GET /auth/profile} without \texttt{Authorization: Bearer ...}. & Request rejected. \\
\hline
AUTH-04 & Refresh token success (\texttt{POST /auth/refresh\_token}) & Send \texttt{refresh} value in the body containing a valid refresh token. & Response contains \texttt{access\_token} and a rotated \texttt{refresh\_token}. \\
\hline
AUTH-05 & Role enforcement & Call admin/relief-only routes (for example \texttt{GET /reports}) using a citizen JWT. & Request rejected due to insufficient role. \\
\hline
AUTH-06 & Admin user listing & Call \texttt{GET /auth/all} with an admin JWT. & Returns list of users with profile fields. \\
\hline
\end{tabularx}
\end{table}

\section{Report and Evidence Scenarios}

\begin{table}[H]
\centering
\caption{Report workflow and evidence handling validation scenarios.}
\label{tab:appendix-report-test-cases}
\renewcommand{\arraystretch}{1.15}
\scriptsize
\setlength{\tabcolsep}{3pt}
\begin{tabularx}{\textwidth}{|p{2.0cm}|p{3.8cm}|X|p{3.2cm}|}
\hline
\textbf{ID} & \textbf{Feature} & \textbf{Input / Steps (summary)} & \textbf{Expected Result} \\
\hline
REP-01 & Create report with evidence (\texttt{POST /reports/create}) & Submit required fields (\texttt{description}, \texttt{province}, \texttt{ward}, \texttt{addressLine}) with \texttt{category[]} and 1--3 image files. & Response indicates success and returns \texttt{reportId}, \texttt{evidences} metadata, and \texttt{images[]} URLs. \\
\hline
REP-02 & Create report without evidence & Submit report payload with no files. & Report created; \texttt{evidences} and \texttt{images} are empty arrays. \\
\hline
REP-03 & Citizen fetches own reports (\texttt{GET /reports/user}) & Call as citizen with valid JWT. & Returns list of reports created by the current user. \\
\hline
REP-04 & Admin/relief fetches all reports (\texttt{GET /reports}) & Call as admin or relief user. & Returns list of report records for review. \\
\hline
REP-05 & Update report by owner (\texttt{PUT /reports/update/:id}) & Update an existing report id owned by the current user; optionally attach additional evidence files. & Returns success message and updated evidence list and image URLs. \\
\hline
REP-06 & Delete report by owner (\texttt{DELETE /reports/:id}) & Delete an existing report id owned by the current user. & Returns success confirmation; subsequent fetch does not include the deleted report. \\
\hline
\end{tabularx}
\end{table}

\section{Tracking Scenarios (Socket.IO)}

\begin{table}[H]
\centering
\caption{Socket.IO live tracking validation scenarios.}
\label{tab:appendix-tracking-test-cases}
\renewcommand{\arraystretch}{1.15}
\scriptsize
\setlength{\tabcolsep}{3pt}
\begin{tabularx}{\textwidth}{|p{2.0cm}|p{3.8cm}|X|p{3.2cm}|}
\hline
\textbf{ID} & \textbf{Event} & \textbf{Input / Steps (summary)} & \textbf{Expected Result} \\
\hline
TRK-01 & Broadcast valid location (\texttt{send-location}) & Connect a client and emit \texttt{send-location} with valid latitude/longitude ranges. & Server broadcasts \texttt{receive-location} containing the sender id and payload. \\
\hline
TRK-02 & Reject invalid payload & Emit \texttt{send-location} with missing or out-of-range coordinates. & Client receives \texttt{location-error} and location is not broadcast. \\
\hline
TRK-03 & Disconnect notification & Disconnect a client socket. & Server broadcasts \texttt{user-disconnected} with the socket id. \\
\hline
\end{tabularx}
\end{table}

\chapter{Runtime Configuration Summary}

This chapter summarizes the environment variables used to deploy VietFlood with Docker Compose. Values shown here are variable names and example formats only. Secret values are not included.

\begin{table}[H]
\centering
\caption{Environment variable groups used for VietFlood deployment with Docker Compose.}
\label{tab:deployment-config-groups}
\scriptsize
\setlength{\tabcolsep}{4pt}
\renewcommand{\arraystretch}{1.1}
\begin{tabularx}{\textwidth}{|p{2.6cm}|>{\raggedright\arraybackslash}X|>{\raggedright\arraybackslash}p{4.2cm}|}
\hline
\textbf{Group} & \textbf{Example Variables} & \textbf{Notes} \\
\hline
Gateway & \texttt{API\_GATEWAY\_PORT}, \texttt{CORS\_ORIGINS} & HTTP port and allowed web origins for the gateway service. \\
\hline
Runtime & \texttt{NODE\_ENV} & Runtime mode used by the Node.js applications (for example, \texttt{production}). \\
\hline
Database & \texttt{DATABASE\_URL} & PostgreSQL connection string for the deployment database. \\
\hline
Message broker & \texttt{RABBITMQ\_DEFAULT\_USER}, \texttt{RABBITMQ\_DEFAULT\_PASS}, \texttt{RABBITMQ\_URL} & RabbitMQ credentials and connection URL used by the backend services. \\
\hline
Cache & \texttt{REDIS\_HOST}, \texttt{REDIS\_PORT}, \texttt{REDIS\_PASSWORD}, \texttt{REDIS\_DB} & Redis connection details used for caching (host, port, password, database index). \\
\hline
Token signing & \texttt{JWT\_SECRET}, \texttt{REFRESH\_SECRET} & Secrets used to sign access tokens and refresh tokens. \\
\hline
Admin seed & \texttt{ADMIN\_EMAIL}, \texttt{ADMIN\_PASSWORD} & Initial administrator credentials used during seed/startup logic. \\
\hline
Cloudinary & \texttt{CLOUDINARY\_CLOUD\_NAME}, \texttt{CLOUDINARY\_API\_KEY}, \texttt{CLOUDINARY\_API\_SECRET} & Cloudinary credentials used for evidence upload and deletion. \\
\hline
\end{tabularx}
\end{table}

\section{Operational Notes}

\begin{itemize}
    \item Use masked values for secrets in the thesis and replace them with actual values in deployment.
    \item Store credentials in CI/CD secrets or environment configuration, not in source control.
    \item Change default RabbitMQ credentials before production use.
    \item Set JWT secrets as cryptographically secure strings.
\end{itemize}

\endgroup
